%% The context category of the registered arrow-only STLC.

\subsection{Simply typed terms as a cartesian closed category}
\label{sec:p10-cartesian-lambda}

Use exactly the types and terms of the registered calculus
\hyperref[sheet:stlc]{\(\mathsf{STLC}\)}: types are generated from
\(\mathsf{Base}\) by \(\alpha\mathbin{\to_{\mathrm{ty}}}\beta\),
and terms are variables, typed abstractions, and applications.
The products and terminal object below are constructions on
contexts.  A tuple of terms is a metalinguistic list.
No product or unit constructor is added to this term language.

\begin{hbdefinition}{def:p10-stlc-beta-eta}{Typed \(\beta\eta\) equality and simultaneous substitution}
Write \(\Theta\vdash_{\beta\eta}M=N:\alpha\) for the least
equivalence relation on terms of type \(\alpha\) in the finite
typing context \(\Theta\) that contains consistent bound
\(\alpha\)-renaming, preserves application and typed abstraction,
and contains the following typed equation schemas:
\[
 (\lambda x{:}\alpha.M)N = M[N/x]:\beta,
 \qquad
 \lambda x{:}\alpha.Mx=M:
                  \alpha\mathbin{\to_{\mathrm{ty}}}\beta.
\]
In the first schema,
\(\Theta,x{:}\alpha\derives_{\mathsf{STLC}}M:\beta\) and
\(\Theta\derives_{\mathsf{STLC}}N:\alpha\).
In the second,
\(\Theta\derives_{\mathsf{STLC}}M:
 \alpha\mathbin{\to_{\mathrm{ty}}}\beta\), and \(x\) is fresh
for \(\Theta\) and \(M\).
Abstraction congruence carries an equation in
\(\Theta,x{:}\alpha\) to an equation in \(\Theta\).
These schemas define an equational metatheory of the registered
syntax; the reduction relation on the calculus sheet is
\(\beta\)-reduction.

A typed simultaneous substitution \(\sigma:\Theta\to\Psi\),
where \(\Psi=(y_1{:}\beta_1,\ldots,y_n{:}\beta_n)\), is a tuple
of terms \(\Theta\derives_{\mathsf{STLC}}\sigma_j:\beta_j\).
For a term \(M\) in \(\Psi\), the notation \(M[\sigma]\)
replaces its free \(y_j\)'s simultaneously by \(\sigma_j\).
Inserted terms are not substituted into again during this
operation.  The operation is capture-avoiding, following
Definition~\ref{def:p07-alpha-substitution}, with binder types
retained.  Under an abstraction, first rename its bound variable
to be fresh for the finitely many terms and contexts involved.
\end{hbdefinition}

\begin{hbtheorem}{thm:p10-stlc-equational-substitution}{proved=annex:p10.lambda.1}{Typing and equality under substitution}
If \(\Psi\derives_{\mathsf{STLC}}M:\alpha\) and
\(\sigma:\Theta\to\Psi\), then
\(\Theta\derives_{\mathsf{STLC}}M[\sigma]:\alpha\).
Identity substitutions act as identities, and for
\(\tau:\Omega\to\Theta\),
\[
 (M[\sigma])[\tau]=_{\alpha}M[\sigma[\tau]],
 \qquad
 \sigma[\tau]=(\sigma_1[\tau],\ldots,\sigma_n[\tau]).
\]
Here \(=_{\alpha}\) denotes bound \(\alpha\)-equivalence.
Both substitution into a typed \(\beta\eta\) equation and
componentwise replacement of a substitution by a
\(\beta\eta\)-equal substitution preserve typed \(\beta\eta\)
equality.  Weakening by unused declarations and consistent
renaming of declared variables preserve typing and equality.
\end{hbtheorem}

\begin{hbdefinition}{def:p10-stlc-context-category}{The category of typed contexts}
Let \(\mathcal S_{\beta\eta}\) have as objects finite type lists
\(\Gamma=(\alpha_1,\ldots,\alpha_m)\), including the empty list
\(()\).  Present \(\Gamma\) by a typing context
\(\Theta_\Gamma=(x_1{:}\alpha_1,\ldots,x_m{:}\alpha_m)\)
with distinct names.  A fixed enumeration of the variable supply
can choose these names once and for all.
An arrow from \(\Gamma\) to
\(\Delta=(\beta_1,\ldots,\beta_n)\) is a componentwise
\(\beta\eta\)-equivalence class of tuples \((M_1,\ldots,M_n)\)
with \(\Theta_\Gamma\derives_{\mathsf{STLC}}M_j:\beta_j\).
Changing the declared names consistently by position changes
only the presentation; this renaming of free context variables
is distinct from bound \(\alpha\)-renaming.

The identity of \(\Gamma\) is the tuple \((x_1,\ldots,x_m)\).
If \(M:\Gamma\to\Delta\) and \(N:\Delta\to\Xi\), define
\[
                 N\circ M=(N_j[M_1/y_1,\ldots,M_n/y_n])_j,
\]
where the \(y_j\)'s name the slots of \(\Delta\).
Fresh names may be used for intermediate presentations.
\end{hbdefinition}

\begin{hbtheorem}{thm:p10-stlc-context-ccc}{proved=annex:p10.lambda.2}{Cartesian closed structure on contexts}
The preceding data form a category.
Its terminal object is \(()\); its product
\(\Gamma\times\Delta\) is the concatenation \(\Gamma\cdot\Delta\),
with block projections and concatenation of tuples as pairing.

For \(\Gamma=(\alpha_1,\ldots,\alpha_m)\), define the
right-associated type
\[
 c_\Gamma(\beta)=
 \alpha_1\mathbin{\to_{\mathrm{ty}}}\cdots
 \mathbin{\to_{\mathrm{ty}}}\alpha_m
 \mathbin{\to_{\mathrm{ty}}}\beta,
 \qquad c_{()}(\beta)=\beta.
\]
If \(\Delta=(\beta_1,\ldots,\beta_n)\), an exponential is
\[
       \Delta^\Gamma=(c_\Gamma(\beta_1),\ldots,c_\Gamma(\beta_n)).
\]
Its evaluation
\(\Delta^\Gamma\cdot\Gamma\to\Delta\) is represented by
\((f_jx_1\cdots x_m)_{j=1}^n\).
The transpose of \(T:\Xi\cdot\Gamma\to\Delta\) is represented by
\[
       \bigl(\lambda x_1{:}\alpha_1\cdots
                       \lambda x_m{:}\alpha_m.T_j\bigr)_{j=1}^n.
\]
For \(m=0\), application and abstraction over the empty list
do nothing.  For \(n=0\), these tuples are empty.
Evaluation and transposition give inverse bijections
\[
 \mathcal S_{\beta\eta}(\Xi\cdot\Gamma,\Delta)
       \cong
 \mathcal S_{\beta\eta}(\Xi,\Delta^\Gamma),
\]
natural in \(\Xi\).  Hence \(\mathcal S_{\beta\eta}\) is
cartesian closed in the sense of Definition~\ref{def:p10-ccc}.
\end{hbtheorem}

\begin{hbdefinition}{def:p10-stlc-ccc-interpretation}{Interpretation of arrow-only STLC in a CCC}
Fix a cartesian closed category \(\mathcal C\) with chosen
terminal object \(1\), binary products, and exponentials, and
assign an object to every base type.  Interpret types by
\[
 \llbracket\alpha\mathbin{\to_{\mathrm{ty}}}\beta\rrbracket
       =\llbracket\beta\rrbracket^{\llbracket\alpha\rrbracket}.
\]
Interpret context lists recursively by
\(\llbracket()\rrbracket=1\) and
\(\llbracket\Gamma\cdot(\alpha)\rrbracket
 =\llbracket\Gamma\rrbracket\times\llbracket\alpha\rrbracket\).
Thus a singleton context is interpreted as
\(1\times\llbracket\alpha\rrbracket\); the construction does
not require this object to be literally
\(\llbracket\alpha\rrbracket\).
A judgment \(\Theta_\Gamma\derives_{\mathsf{STLC}}M:\alpha\)
denotes an arrow
\(\llbracket M\rrbracket:
 \llbracket\Gamma\rrbracket\to\llbracket\alpha\rrbracket\).
Variables denote the corresponding product projections,
application denotes
\[
 \llbracket MN\rrbracket
      =\operatorname{ev}\circ
           \langle\llbracket M\rrbracket,\llbracket N\rrbracket\rangle,
\]
and abstraction denotes
\(\llbracket\lambda x{:}\alpha.M\rrbracket
  =\Lambda(\llbracket M\rrbracket)\), where \(\Lambda\) is
exponential transposition.
The interpretation of a typed substitution is the unique
arrow into the interpreted target context with the specified
term components.
\end{hbdefinition}

\begin{hbtheorem}{thm:p10-stlc-ccc-completeness}{proved=annex:p10.lambda.3}{Substitution, soundness, and completeness for CCC equality}
Interpretation respects bound \(\alpha\)-renaming and satisfies
\[
             \llbracket M[\sigma]\rrbracket
                 =\llbracket M\rrbracket\circ\llbracket\sigma\rrbracket.
\]
For terms \(M,N\) of the same type \(\alpha\) in the same
context \(\Theta\), the following are equivalent:
\[
 \Theta\vdash_{\beta\eta}M=N:\alpha
 \quad\Longleftrightarrow\quad
 \llbracket M\rrbracket=\llbracket N\rrbracket
 \text{ in every CCC interpretation of the base types.}
\]
In particular, the context category itself detects every
failure of typed \(\beta\eta\) equality.
\end{hbtheorem}

Pitts derives the interpretation of function types and the
categorical \(\beta\) and \(\eta\) equations in
\cite[\S3.3, pp.~22--24, equations (3.19)--(3.22)]{Pitts1995Draft}.
His classifying-context construction uses tuples and
simultaneous substitution, with finite products given by
concatenation
\cite[\S4.2, pp.~30--33, Proposition~4.2.1]{Pitts1995Draft}.
The curried-list exponential is also given in
\cite[\S4.4, p.~38]{Pitts1995Draft}.
Its use for this handbook's arrow-only syntax is proved in the annex.
All page locators here refer to the draft dated 18 May 1995.
\index{lambda calculus!cartesian closed interpretation}
\index{cartesian closed category!typed context category}
